% test-020.tex -- regresion de \spcoord
%
% Compilar:  lualatex-dev test-020.tex
% Validar:   show-pdf-tags --xml test-020.pdf | rnv-wrapp
%            verapdf --flavour ua2 --format text test-020.pdf
%
% Test de regresion dedicado a \spcoord, que tuvo varios fallos que se
% pisaban entre si:
%   - la altura de los parentesis (una fraccion en cualquier posicion de la
%     abscisa o de la ordenada, o un \spmQ, los agranda; dentro de un
%     subindice o superindice no);
%   - MathCAT leia (1,2) como intervalo si despues venia = o \in (se
%     envuelve en \mathord);
%   - un unico nodo con raiz, potencia o subindice perdia su arg (coma
%     perdida) si ademas llevaba una fraccion.
% Los "doc:" son lecturas reales de NVDA con MathCAT (2026-10-07), salvo los
% marcados PENDIENTE. "es no es igual a" y "es miembro de a" son lecturas
% propias de MathCAT, no del paquete. El tamano se comprueba en el XML
% (maxsize/minsize del <mo>): 11.907pt en texto, 17.861pt con dfrac o en display.
\DocumentMetadata
  {
    lang=es-CL,
    pdfversion=2.0,
    pdfstandard=ua-2,
    tagging=on,
    tagging-setup={math/setup={mathml-SE}},
  }
\documentclass{article}
\usepackage[spanish]{babel}
\usepackage{unicode-math}
\usepackage{spintent}
\usepackage{hyperref}
\hypersetup
  {
    pdfdisplaydoctitle = true,
    pdftitle           = {test-020: regresion de spcoord},
  }
\begin{document}

\section*{A. Varias coordenadas en una fórmula}

% doc: «1 coma 2 es igual a 3 coma 4»
1. Dos coordenadas con igual: $\spcoord{1,2} = \spcoord{3,4}$.

% doc: «1 coma 2 más 3 coma 4»
2. Dos coordenadas con suma: $\spcoord{1,2} + \spcoord{3,4}$.

% doc: «x es igual a 3 coma 4»
3. Una coordenada tras el igual: $x = \spcoord{3,4}$.

% doc: «1 coma 2 es igual a 3 coma 4 es igual a 5 coma 6»
4. Tres coordenadas encadenadas: $\spcoord{1,2} = \spcoord{3,4} = \spcoord{5,6}$.

% doc: «punto A 1 coma 2 es igual a punto B 3 coma 4»
5. Dos puntos con coordenadas: $\spPoint{A}(1,2) = \spPoint{B}(3,4)$.

\section*{B. Lo que viene después del paréntesis (MathCAT leía el par como un intervalo)}

% doc: «1 coma 2 es igual a 5»
6. Igual y un número: $\spcoord{1,2} = 5$.

% doc: «1 coma 2 es miembro de a»
7. Pertenencia: $\spcoord{1,2} \in A$.

% doc: «1 coma 2 es no es igual a 3 coma 4»
8. Distinto: $\spcoord{1,2} \neq \spcoord{3,4}$.

% doc: «1 coma 2 es menor que 3 coma 4»
9. Menor: $\spcoord{1,2} < \spcoord{3,4}$.

% doc: «1 coma 2 3»
10. Sin operador, producto implícito: $\spcoord{1,2}\,3$.

% doc: «1 punto y coma 2 es igual a 5»
11. Con punto y coma: $\spcoord{1;2} = 5$.

% doc: «x es igual a 1 coma 2 es igual a 5»
12. Una variable antes y un igual después: $x = \spcoord{1,2} = 5$.

\section*{C. Tamaño de los paréntesis (mirar maxsize en el XML)}

% doc: «1 mitad coma 3»
13. Fracción al comienzo (11.907pt): $\spcoord{\frac{1}{2}, 3}$.

% doc: «1 más 1 mitad coma 3»
14. Fracción después de un término, en la abscisa (11.907pt): $\spcoord{1+\frac{1}{2}, 3}$.

% doc: «3 coma 1 más 1 mitad»
15. Fracción después de un término, en la ordenada (11.907pt): $\spcoord{3, 1+\frac{1}{2}}$.

% doc: «1 más 1 mitad coma 3»
16. Con dfrac después de un término (17.861pt): $\spcoord{1+\dfrac{1}{2}, 3}$.

% doc: «el raíz cuadrada de 1 mitad coma 1»
17. Fracción dentro de una raíz (11.907pt): $\spcoord{\sqrt{\frac{1}{2}}, 1}$.

% doc: «x subíndice 1 mitad coma 1»
18. Fracción dentro de un subíndice, sin tamaño: $\spcoord{x_{\frac{1}{2}}, 1}$.

% doc: «1 y la fracción con numerador 1 y denominador 2 coma 3»
19. Mixto en la abscisa (11.907pt): $\spcoord{\spmQ{1}{1}{2}, 3}$.

% doc: «3 coma 1 y la fracción con numerador 1 y denominador 2»
20. Mixto en la ordenada (11.907pt): $\spcoord{3, \spmQ{1}{1}{2}}$.

% doc: «1 y la fracción con numerador 1 y denominador 2»
21. Un mixto solo, el control: $\spmQ{1}{1}{2}$.

\section*{D. Un único nodo con raíz, potencia o subíndice (necesita el separador interno)}

% doc: «el raíz cuadrada de 2 coma 1»
22. Una raíz en la abscisa: $\spcoord{\sqrt{2}, 1}$.

% doc: «x elevado al cuadrado coma 1»
23. Una potencia en la abscisa: $\spcoord{x^{2}, 1}$.

% doc: «1 coma x subíndice 1»
24. Un subíndice en la ordenada: $\spcoord{1, x_{1}}$.

% doc: «3 coma 1 mitad»
25. Con dfrac en la ordenada (17.861pt): $\spcoord{3, \dfrac{1}{2}}$.

\section*{E. Fórmula destacada}

% doc: «1 más 1 mitad coma 3»
26. Fracción después de un término, en display (17.861pt):
\[ \spcoord{1+\frac{1}{2}, 3} \].

\end{document}
